%-----------------------------------------------------------------
% Rodrigo Chang, 2019
% Banco de Guatemala
% Programa de Estudios Superiores 2019-2020
% Diapositivas del curso "Programación I"
%
% Descargo de responsabilidad.
%   El contenido de este documento es responsabilidad exclusiva de
% su(s) autores. Cualquier opinión expresada no necesariamente
% representa la postura institucional del Banco de Guatemala. 
% 
% Contenido.
%   El contenido de estas diapositivas está construido a partir de 
% modificaciones al material provisto por Ana Bell, Eric Grimson y 
% John Guttag para el curso 6.0001 Introduction to Computer Science 
% and Programming in Python, otoño 2016, Massachusetts Institute of 
% Technology: MIT OpenCourseWare, https://ocw.mit.edu
%
% Licencia: 
% Creative Commons BY-NC-SA
% https://creativecommons.org/licenses/by-nc-sa/4.0/.
%
% Formato. 
%   El formato utilizado para la elaboración de este documento 
% ha sido modificado a partir del repositorio: 
%       https://github.com/rinnert/statlearning-python-pwf
%
% Kurt Rinnert, 2019
%   Copying and distribution of this file, with or without 
% modification, are permitted in any medium without royalty provided 
% the copyright notice and this notice are preserved. This file is 
% offered as-is, without any warranty.
%
%-----------------------------------------------------------------
\documentclass[mathserif]{beamer}
% Para configurar las diapositivas en relación 16:9, 
% agregar la opción aspectratio=169

% Título de la presentación y tópicos a tratar
\usepackage{arev, arevmath}
\usepackage[scaled]{cabin}
%\usepackage[T1]{fontenc}
\usepackage[super]{nth}
\usepackage{pifont}
\usepackage{wasysym}
\usepackage{tabularx}
\usepackage{array}
\usepackage{booktabs}
\usepackage{boldline}
\usepackage{colortbl}
\usepackage{amsmath}
\usepackage{bm}
\usepackage{tcolorbox}
\usepackage{adjustbox}
\usepackage{minibox}
\usepackage{makecell}
\usepackage{adjustbox}
\usepackage{textcomp}
\usepackage[absolute,overlay]{textpos}
\setlength{\TPHorizModule}{1mm}%
\setlength{\TPVertModule}{1mm}%
\tcbuselibrary{skins}
\usepackage{graphicx}
\usepackage[spanish, es-tabla]{babel}
\usepackage[utf8]{inputenc}
\usepackage{csquotes}

\usepackage{comment}
\usepackage{svg}

\usepackage{verbatim}
\usepackage{pdfpages}

% Columnas múltiples
\usepackage{multicol}

% -------------------------------------------
%	LISTINGS
% -------------------------------------------
% Paquetes necesarios para los códigos
\usepackage{listings}
\usepackage{color}
% Configuración de listings (pendiente)
%\input{conf/listingStyles}
% Redefinición de titulos de caption
\renewcommand\lstlistingname{\textbf{Código}}
\renewcommand\lstlistlistingname{Códigos}

% Definición de colores para Python (Jupyter notebooks)
\definecolor{python_keyword_color}{RGB}{0, 128, 0}
\definecolor{python_string_color}{RGB}{186, 34, 34}
\definecolor{python_comment_color}{RGB}{1, 121, 121}
\definecolor{frame_numbers_color}{rgb}{0.5,0.5,0.5}
\definecolor{python_identifier_color}{RGB}{0,0,0}
\definecolor{jupyter_background}{RGB}{247, 247, 247}
\definecolor{jupyter_frame}{RGB}{207, 207, 207}

% Definicion de estilo para Python
\lstdefinestyle{custom_python}{
	language=python,			% the language of the code
	backgroundcolor=\color{jupyter_background}, 
	% the size of the fonts that are used for the code
	basicstyle=\linespread{1}\small\ttfamily,	
	% Configuración de colores
	commentstyle=\color{python_comment_color},	% comment style
	identifierstyle=\color{python_identifier_color},
	keywordstyle=\mdseries\color{python_keyword_color},	% keyword style
	numberstyle=\tiny\color{frame_numbers_color},	
	stringstyle=\color{python_string_color}, % string literal style 
	captionpos=t,				% sets the caption-position to bottom
	breakatwhitespace=false,	% sets if automatic breaks should only happen at whitespace
	breaklines=true,			% sets automatic line breaking
	morekeywords={},			% if you want to add more keywords to the set
	deletekeywords={eps},		% delete keywords from the given language
	escapeinside={\%*}{*)},		% if you want to add LaTeX within your code
	framesep=3pt,
	keepspaces=true,			% keeps spaces in text,       
	numbers=none, 	% (none, left, right)
	numbersep=5pt,	% how far the line-numbers are from the code
	% if not set, the frame-color may be changed on line-breaks within not-black text (e.g. comments (green here))
	rulecolor=\color{jupyter_frame}, 
	% show spaces everywhere adding particular underscores; it overrides 'showstringspaces'
	showspaces=false,                
	% underline spaces within strings only
	showstringspaces=false,          
	% show tabs within strings adding particular underscores
	showtabs=false,                  
	stepnumber=0,	% the step between two line-numbers      
	% sets default tabsize to 2 spaces
	tabsize=2,                       
	% lets you use non-ASCII characters; for 8-bits encodings only, does not work with UTF-8
	extendedchars=true,      
	% Support additional characters            
	literate=
  {á}{{\'a}}1 {é}{{\'e}}1 {í}{{\'i}}1 {ó}{{\'o}}1 {ú}{{\'u}}1
  {Á}{{\'A}}1 {É}{{\'E}}1 {Í}{{\'I}}1 {Ó}{{\'O}}1 {Ú}{{\'U}}1
  {à}{{\`a}}1 {è}{{\`e}}1 {ì}{{\`i}}1 {ò}{{\`o}}1 {ù}{{\`u}}1
  {À}{{\`A}}1 {È}{{\'E}}1 {Ì}{{\`I}}1 {Ò}{{\`O}}1 {Ù}{{\`U}}1
  {ä}{{\"a}}1 {ë}{{\"e}}1 {ï}{{\"i}}1 {ö}{{\"o}}1 {ü}{{\"u}}1
  {Ä}{{\"A}}1 {Ë}{{\"E}}1 {Ï}{{\"I}}1 {Ö}{{\"O}}1 {Ü}{{\"U}}1
  {â}{{\^a}}1 {ê}{{\^e}}1 {î}{{\^i}}1 {ô}{{\^o}}1 {û}{{\^u}}1
  {Â}{{\^A}}1 {Ê}{{\^E}}1 {Î}{{\^I}}1 {Ô}{{\^O}}1 {Û}{{\^U}}1
  {œ}{{\oe}}1 {Œ}{{\OE}}1 {æ}{{\ae}}1 {Æ}{{\AE}}1 {ß}{{\ss}}1
  {ű}{{\H{u}}}1 {Ű}{{\H{U}}}1 {ő}{{\H{o}}}1 {Ő}{{\H{O}}}1
  {ç}{{\c c}}1 {Ç}{{\c C}}1 {ø}{{\o}}1 {å}{{\r a}}1 {Å}{{\r A}}1
  {€}{{\EUR}}1 {£}{{\pounds}}1,   
}

% Configurar el listing para el documento actual
\lstset{
	style = custom_python, 
	frame = single
}

% -------------------------------------------
%	OTRAS CONFIGURACIONES
% -------------------------------------------
\makeatletter
\newcommand{\antsize}{\@setfontsize{\antsize}{4pt}{4pt}}
\makeatother
\newcommand{\at}{\makeatletter @\makeatother}

\newcommand{\cmark}{\ding{51}}%
\newcommand{\bottomline}[1]{\vskip0pt plus 1fill{\alert{#1}\phantom{g}\vskip 1.0mm}}

% Rutas para las imágenes
%\graphicspath{{figs/}{figs/common/}}

% Separación de filas en las tablas más amplia
\renewcommand{\arraystretch}{1.2}


% -------------------------------------------
%	ESTILO DE LA PRESENTACIÓN
% -------------------------------------------

\mode<presentation>%
{
	\usetheme{default}
	%\usetheme[width=2.5cm]{PaloAlto}
	\usecolortheme{dove}
	\useoutertheme{infolines}
	% oder auch nicht

	% BANGUAT Colors
	\definecolor{BANGUATblue}{RGB}{37,86,162} % 0x255682
	\definecolor{BANGUATorange}{RGB}{255,130,0} % 0xff8200
	\definecolor{BANGUATgreen}{RGB}{0,100,0}
	\definecolor{BANGUATdark}{RGB}{80,80,80} % 0x505050
	\definecolor{BANGUATlight}{RGB}{120,120,120}
	\definecolor{BANGUATbrown}{RGB}{178,91,0}

	\definecolor{codebg}{rgb}{0.95,0.95,0.95}

	% Color theme
	\setbeamercolor{alerted text}{fg=BANGUATorange}
	\setbeamercolor{frametitle}{fg=BANGUATblue}
	\setbeamercolor{title}{fg=BANGUATblue}
	\setbeamercolor{subtitle}{fg=BANGUATorange}
	\setbeamercolor{normal text}{fg=BANGUATdark}
	\setbeamercolor{author in foot}{fg=BANGUATblue}
	\setbeamercolor{item}{fg=BANGUATblue}
	\setbeamercolor{footline}{fg=BANGUATblue}
	%\setbeamercolor{item projected}{bg=BANGUATorange}
	%\setbeamercolor{item projected}{fg=white}

	\setbeamertemplate{headline}
	{}
	\setbeamertemplate{frametitle}
	{
		%\textbf{{\insertframetitle\phantom{g}}}\\
		%\textbf{\insertframetitle\phantom{g}}\\
		\textbf{\underline{\insertframetitle\phantom{g}}}\\
		%\textbf{\underline{\insertframetitle}}\\
		\vskip 1.0mm
		%{\color{UOLgold}\hrule height 2pt}
		%\par
	}
	\addtobeamertemplate{frametitle}{}{\vspace{-1em}}
	\setbeamertemplate{footline}{
		{%
			\textbf{ \hskip 3.0mm\insertshorttitle\phantom{.}---\phantom{.}\insertshortinstitute\hfill\insertframenumber\,/\,\inserttotalframenumber\hskip 3.0mm} 
		}
	}

	\setbeamertemplate{navigation symbols}{}%remove navigation symbols
	\setbeamertemplate{itemize items}[circle]
	\setbeamertemplate{enumerate items}[fg=BANGUATblue]
	\setbeamercolor{itemize items}{fg=BANGUATblue}
	\setbeamercolor{sidebar}{bg=BANGUATblue}
	\setbeamercolor{title in sidebar}{fg=BANGUATorange}
	\setbeamercolor{author in sidebar}{fg=BANGUATorange}
	\setbeamercolor{section in sidebar}{fg=BANGUATorange}
}



% -------------------------------------------
%	CONFIGURACIÓN DE DIAPOSITIVA DE TÍTULO
% -------------------------------------------

% Título del curso
\title[Programación I]{\textbf{%
Introducción a las ciencias de la computación\\%
{\it y programación en Python}}\\%
%\href{https://ocw.mit.edu/courses/electrical-engineering-and-computer-science/6-0001-introduction-to-computer-science-and-programming-in-python-fall-2016/}{\tiny\it Basado en el curso ``Introduction to Computer Science and Programming in Python'' por Ana Bell, Eric Grimson, y John Guttag.}%
\vspace{3em}%
}
%\vspace{-2.5cm}{}

% Autor y correo electrónico
\author{\href{mailto:?to=Rodrigo Chang <rrcp@banguat.gob.gt>&subject=PES Programación I}{Rodrigo Chang}}

% Información del instituto
\institute[\href{https://www.banguat.gob.gt/}{Banco de Guatemala}]%
{\color{BANGUATblue} \bfseries \href{https://www.banguat.gob.gt/}{Banco de Guatemala}\\%
\vspace{3mm}%
\includegraphics[width=0.13\textwidth]{figs/common/logo-banguat.png}}

% Sin fecha en particular
\date{}

% Información de licencias y copyright
\titlegraphic{
	 %\texorpdfstring{\vspace{-3cm}}{}
	 \vspace{-1.5cm}
	 \begin{minipage}[b][1cm][b]{0.25\textwidth}%
	    \color{BANGUATlight}\antsize
		%Copyright \textcopyright~2019\\
		\href{mailto:?to=Rodrigo Chang <rrcp@banguat.gob.gt>&subject=PES Programación I}{Rodrigo Chang <rrcp@banguat.gob.gt>}\\
		Este material está construido a partir de modificaciones al material provisto por Ana Bell, Eric Grimson y John Guttag para el curso \textit{6.0001 Introduction to Computer Science and Programming in Python}, otoño 2016, Massachusetts Institute of Technology: MIT OpenCourseWare, \href{https://ocw.mit.edu}{https://ocw.mit.edu}. Licencia: \href{https://creativecommons.org/licenses/by-nc-sa/4.0/}{Creative Commons BY-NC-SA}.
		
		%Copying and distribution of this file, with or without modification, are permitted in any medium without royalty provided the copyright notice and this notice are preserved.  This file is offered as-is, without any warranty.


		%Some of the figures in this presentation are taken from ``An Introduction to Statistical Learning, with applications in R''  (Springer, 2013) with permission from the authors: G. James, D. Witten,  T. Hastie and R. Tibshirani 
	 \end{minipage}%
	 \hspace{9cm}
}

% Restar uno al contador
\addtocounter{framenumber}{-1}

% Cajas de colores
\newcommand{\tabboxset}{\tcbset{enhanced, nobeforeafter, boxrule=0pt, boxsep=0pt, colback=codebg, colframe=codebg, coltext=BANGUATdark, rounded corners, arc=4pt, fonttitle={\bfseries\tiny}}}
\newcommand{\codeboxset}{\tcbset{enhanced, nobeforeafter, boxrule=0pt, boxsep=0pt, colback=codebg, colframe=codebg, coltext=BANGUATdark, rounded corners, arc=4pt, fonttitle={\bfseries\tiny}}}

% Comandos útiles
\newcommand{\orange}{\color{BANGUATorange}}
\newcommand{\blue}{\color{BANGUATblue}}
\newcommand{\dark}{\color{BANGUATdark}}
\newcommand{\R}{\mathbb{R}}
\newcommand{\dat}[1]{{\footnotesize\tt\orange #1}}
\newcommand{\mdat}[1]{\text{\dat{#1}}}
\newcommand{\val}[1]{{\footnotesize\tt\blue #1}}
\newcommand{\mval}[1]{\text{\val{#1}}}
\newcommand{\e}[1]{\emph{#1}}
\newcommand{\bh}{\hat{\beta}}
\newcommand{\h}{\hat}
\newcommand{\bl}{\bottomline}


\makeatletter
\newcommand{\includegraphicsdpi}[3]{%
	\pdfimageresolution=#1%
	\includegraphics[#2]{#3}%
	\pdfimageresolution=72%
}

\newcommand{\Quote}[2]{%
	\begin{center} 
		\begin{minipage}{0.7\textwidth} 
			\hrule
			\vskip 3mm
			\emph{{\color{BANGUATblue} ``#1''}\\
			{\color{BANGUATorange} --- #2}}
			\vskip 3mm
			\hrule
			\vskip 2mm
		\end{minipage}
	\end{center}
}

% Texto centrado en la diapositiva, quizás para el abstract o para indicar abajo de un quote
\newenvironment{blurb}%
	{\begin{center}\begin{minipage}{0.6\textwidth}\footnotesize}
	{\end{minipage}\end{center}}

% Environment para códigos o algoritmos. Parecido a un bottomline, se puede utilizar para dar definiciones o enfatizar en algún concepto
\newenvironment{cpage}%
	{\begin{center}\begin{minipage}{0.75\textwidth}}
	{\end{minipage}\end{center}}

% Bloque (similar a un pop-up) de alerta para indicar resultados en tablas, definiciones, fórmulas, etc.
% Primer argumento es el tamaño 
% Segundo argumento es el título del bloque
\newenvironment{popblock}[2]%
	{\begin{center}\begin{minipage}{#1}\footnotesize
		\begin{tcolorbox}[colframe=codebg, colback=white, colupper=BANGUATdark, title={\normalsize\bfseries\blue #2}]}
	{\end{tcolorbox}\end{minipage}\end{center}}
\makeatother

% Incluir diapositivas en PDF
% 1 - Título de la diapositiva
% 2 - páginas del PDF a incluir
% 3 - ruta hacia el archivo PDF
\newcommand{\includeslides}[3]{%
{%
\setbeamercolor{background canvas}{bg=}%
\addtocounter{framenumber}{1}%
\includepdf[pages={#1}, link=true, pagecommand={%
\Large\blue\bfseries\underline{%
#3\phantom{g}}%
\vspace{-1em}%
\vskip 1.0mm}]{#2}%
}}
\subtitle{%
  \bfseries{Tuplas y listas: mutabilidad y clonado}\\%
  %{\tiny\it Manejo de strings, \textit{guess and check}, métodos de aproximación}\\%
}

% Inicio del documento
\begin{document}
\frame[plain]{
	%\vskip 1.0mm
	\titlepage
	%\vskip 1.0mm
}
%-----------------------------------------------------------------
% Empieza el contenido de la presentación
%-----------------------------------------------------------------

\begin{frame}{Abstract}
	\Quote{La simplicidad es un prerrequisito para la confiabilidad}{Edsger W. Dijkstra}
    
	\begin{blurb}
		Veremos dos tipos nuevos {\blue compuestos} en Python.
		
		Hablaremos de los conceptos de \textit{alias}, mutabilidad y clonado.
		
	\end{blurb}
\end{frame}

\section{Tuplas}

\begin{frame}[fragile]{Tuplas}
    \begin{itemize}
        \item Secuencia {\blue ordenada} de elementos, que pueden ser de diferentes {\blue tipos}.
        \item No es posible alterar sus elementos, son {\bfseries\orange inmutables}.
        \item Se representan con paréntesis
    \end{itemize}
\begin{lstlisting}
# Tupla vacia
te = ()
# Tupla de ejemplo
t = (2, "pes", 3.2)

t[0]      # devuelve 2
t + (5,6) # devuelve (2, "pes" 3.2, 5, 6)
t[1:2]    # slice, devuelve ("pes", )
t[1:3]    # slice, devuelve ("pes", 3.2)
len(t)    # devuelve 3

t[1] = 4 # error, inmutable
\end{lstlisting}
\end{frame}

% Para qué sirven las tuplas: cambiar valores entre sí y devolver más de un valor
\includeslides{5}{slides-60001/MIT6_0001F16_Lec5_mod.pdf}{Tuplas}

% Ejemplo de iteración sobre tuplas\textbf{}
\begin{frame}[fragile]{Iteración sobre tuplas}
    \begin{itemize}
        \item Es posible {\blue iterar} sobre los elementos de una tupla:
    \end{itemize}
\begin{lstlisting}
t = (2, "pes", 3.2)

for elem in t: 
    print(elem, "es de tipo", type(elem))
\end{lstlisting}
\end{frame}

% Iteration over tuples with pairs
\begin{frame}[fragile]{Tuplas de tuplas}
\begin{lstlisting}
def get_data(aTuple):
    nums= ()
    words = ()
    for t in aTuple:
        nums= nums+ (t[0],)
        if t[1] not in words:
            words = words + (t[1],)
    
    min_n = min(nums)
    max_n = max(nums)
    unique_words = len(words)
    return (min_n, max_n, unique_words)
    
aTuple = ((1, 'a'), (2, 'b'), (3, 'c'))
print(get_data(aTuple))
\end{lstlisting}
\end{frame}

\section{Listas}
\begin{frame}[fragile]{Listas}
    \begin{itemize}
        \item {\blue\bfseries Secuencia ordenada} de elementos, accesibles a través de un {\blue índice}.
        \item Se denotan utilizando {\blue corchetes}, \lstinline|[ ]|
        \item Los elementos: 
            \begin{itemize}
                \item Son usualmente {\orange homogéneos}
                \item Pero pueden ser de diferentes tipos (práctica poco común, pero válida).
            \end{itemize}
        
        \item Los elementos pueden ser {\orange\bfseries alterados}, por lo que la lista es {\orange\bfseries mutable}.
    \end{itemize}
\end{frame}

\begin{frame}[fragile]{Lista e índices}
\begin{lstlisting}
L = []
L = [2, 'a', 4, [1,2]]

len(L) # evaluates to 4
L[0]   # evaluates to 2
L[2]+1 # evaluates to 5
L[3]   # evaluates to [1,2], another list!
L[4]   # gives an error

i = 2
L[i-1] # evaluates to 'a' since L[1]='a' above
\end{lstlisting}
\end{frame}

% Cambiar elementos
\includeslides{9}{slides-60001/MIT6_0001F16_Lec5_mod.pdf}{Cambiando elementos}

% Iterar sobre una lista
\includeslides{10}{slides-60001/MIT6_0001F16_Lec5_mod.pdf}{Iterando sobre una lista}

% Operaciones sobre listas
\includeslides{11-13}{slides-60001/MIT6_0001F16_Lec5_mod.pdf}{Operaciones sobre listas}

% Convertir listas a strings
\includeslides{14}{slides-60001/MIT6_0001F16_Lec5_mod.pdf}{Convertir listas a cadenas de texto}

% Otras operaciones sobre listas
\begin{frame}[fragile]{Otras operaciones sobre listas}
    \begin{itemize}
        \item Funciones para ordenar: \lstinline|sort()| y \lstinline|sorted()|
        \item Operación de {\blue reversa}: \lstinline|reverse()|
        \item Otras operaciones más: \url{https://docs.python.org/3/tutorial/datastructures.html}
    \end{itemize}
    
    {\blue\bfseries ¡Ojo!} algunas son {\blue funciones} y otras son {\orange métodos}. 
    
\begin{lstlisting}
L = [9,6,0,3]
sorted(L)   # devuelve lista ordenada sin cambiar L
L.sort()    # cambia L = [0, 3, 6, 9]
L.reverse() # cambia L = [9, 6, 3, 0]
\end{lstlisting}
\end{frame}


\section{Mutabilidad y clonado}

\begin{frame}[fragile]{Listas en memoria}
    \begin{itemize}
        \item Las listas son objetos {\blue mutables}.
        \item Se comportan de forma diferente a los tipos inmutables.
        \item Representan {\orange objetos en memoria}. 
        \item Distintos {\orange nombres de variables} pueden apuntar al mismo objeto en memoria.
        \item Cualquier variable que {\blue\bfseries apunte } a un objeto es afectada.
    \end{itemize}
    
    \bottomline{Cuando trabajemos con lista debemos cuidarnos de los efectos secundarios.}
\end{frame}

\begin{frame}[fragile]{Alias para listas}
    \begin{itemize}
        \item \lstinline|hot| es un alias para \lstinline|warm|: al cambiar uno, cambia el otro.
        
        \item La función \lstinline|append()| tiene un {\blue efecto secundario.}

\begin{lstlisting}
a = 1
b = a
print(a)
print(b)

warm = ['red', 'yellow', 'orange']
hot = warm
hot.append('pink')
print(hot)
print(warm)
\end{lstlisting}

        \item \href{http://pythontutor.com/visualize.html#code=a%20%3D%201%0Ab%20%3D%20a%0Aprint%28a%29%0Aprint%28b%29%0A%0Awarm%20%3D%20%5B'red',%20'yellow',%20'orange'%5D%0Ahot%20%3D%20warm%0Ahot.append%28'pink'%29%0Aprint%28hot%29%0Aprint%28warm%29&cumulative=false&curInstr=0&heapPrimitives=nevernest&mode=display&origin=opt-frontend.js&py=3&rawInputLstJSON=%5B%5D&textReferences=false}{Ver en Python Tutor}
        

    \end{itemize}
\end{frame}


\begin{frame}[fragile]{Clonando una lista}
    \begin{itemize}
        \item Para crear una nueva lista podemos {\blue\bfseries copiar cada elemento} utilizando \textit{slicing}. 
        
        \item \href{http://pythontutor.com/visualize.html#code=cool%20%3D%20%5B'blue',%20'green',%20'grey'%5D%0Achill%20%3D%20cool%5B%3A%5D%0Achill.append%28'black'%29%0Aprint%28chill%29%0Aprint%28cool%29&cumulative=false&heapPrimitives=nevernest&mode=edit&origin=opt-frontend.js&py=3&rawInputLstJSON=%5B%5D&textReferences=false}{Ver en Python Tutor}
    \end{itemize}
\begin{lstlisting}
cool = ['blue', 'green', 'grey']
chill = cool[:]
chill.append('black')
print(chill)
print(cool)
\end{lstlisting}
\end{frame}


\begin{frame}[fragile]{Ordenando listas}
    \begin{itemize}
        \item Utilizar \lstinline|sort()| altera la lista, devuelve \lstinline|None|. 
        \item Utilizar \lstinline|sorted()| {\blue no altera} la lista, devuelve la lista ordenada. 

\begin{lstlisting}
warm = ['red', 'yellow', 'orange']
sortedwarm = warm.sort()
print(warm)
print(sortedwarm)

cool = ['grey', 'green', 'blue']
sortedcool = sorted(cool)
print(cool)
print(sortedcool)
\end{lstlisting}

        \item \href{http://pythontutor.com/visualize.html#code=warm%20%3D%20%5B'red',%20'yellow',%20'orange'%5D%0Asortedwarm%20%3D%20warm.sort%28%29%0Aprint%28warm%29%0Aprint%28sortedwarm%29%0A%0Acool%20%3D%20%5B'grey',%20'green',%20'blue'%5D%0Asortedcool%20%3D%20sorted%28cool%29%0Aprint%28cool%29%0Aprint%28sortedcool%29&cumulative=false&heapPrimitives=nevernest&mode=edit&origin=opt-frontend.js&py=3&rawInputLstJSON=%5B%5D&textReferences=false}{Ver en Python Tutor}
    \end{itemize}
\end{frame}

\begin{frame}[fragile]{Listas de listas de listas de ...}
    \begin{itemize}
        \item Es posible {\blue\bfseries agrupar} una lista dentro de otra.
        \item Cuidado con los {\orange\bfseries efectos secundarios}.
    
\begin{lstlisting}
warm = ['yellow', 'orange']
hot = ['red']
brightcolors = [warm]
brightcolors.append(hot)
print(brightcolors)

hot.append('pink')
print(hot)
print(brightcolors)
\end{lstlisting}

        \item \href{http://pythontutor.com/visualize.html#code=warm%20%3D%20%5B'yellow',%20'orange'%5D%0Ahot%20%3D%20%5B'red'%5D%0Abrightcolors%20%3D%20%5Bwarm%5D%0Abrightcolors.append%28hot%29%0Aprint%28brightcolors%29%0A%0Ahot.append%28'pink'%29%0Aprint%28hot%29%0Aprint%28brightcolors%29&cumulative=false&heapPrimitives=nevernest&mode=edit&origin=opt-frontend.js&py=3&rawInputLstJSON=%5B%5D&textReferences=false}{Ver en Python Tutor}
    
    \end{itemize}
\end{frame}

% Mutación e iteración - evitar mutar una lista al iterar sobre ella
\begin{frame}[fragile]{}
    \begin{itemize}
        \item Como {\blue\bfseries buena práctica}, es mejor evitar mutar una lista al iterar sobre ella. Veamos un ejemplo: 
        
        \item \href{http://pythontutor.com/visualize.html#code=def%20remove_dups%28L1,%20L2%29%3A%0A%20%20%20%20for%20e%20in%20L1%3A%0A%20%20%20%20%20%20%20%20if%20e%20in%20L2%3A%0A%20%20%20%20%20%20%20%20%20%20%20%20L1.remove%28e%29%0A%0AL1%20%3D%20%5B1,%202,%203,%204%5D%0AL2%20%3D%20%5B1,%202,%205,%206%5D%0Aremove_dups%28L1,%20L2%29%0Aprint%28L1,%20L2%29&cumulative=false&curInstr=15&heapPrimitives=nevernest&mode=display&origin=opt-frontend.js&py=3&rawInputLstJSON=%5B%5D&textReferences=false}{Ver en Python Tutor}
        
        \item ¿Por qué \lstinline|L1 = [2,3,4]| y no \lstinline|[3,4]|? ¿Cómo se puede corregir este código?
    \end{itemize}
\begin{lstlisting}
def remove_dups(L1, L2):
    for e in L1:
        if e in L2:
            L1.remove(e)

L1 = [1, 2, 3, 4]
L2 = [1, 2, 5, 6]
remove_dups(L1, L2)
print(L1, L2)
\end{lstlisting}
\end{frame}

% Corrección
\begin{frame}[fragile]{}
    \begin{itemize}
        \item Solución.
    \end{itemize}
\begin{lstlisting}
def remove_dups_new(L1, L2):
    L1_copy = L1[:]
    for e in L1_copy:
        if e in L2:
            L1.remove(e)
            
L1 = [1, 2, 3, 4]
L2 = [1, 2, 5, 6]
remove_dups_new(L1, L2)
print(L1, L2)

\end{lstlisting}
\end{frame}

\begin{frame}{}
\centering
{\bfseries\blue\LARGE Ejercicios}
\end{frame}

%-----------------------------------------------------------------
% Termina el contenido de la presentación
%-----------------------------------------------------------------

\begin{comment}

% Slide with two columns, half and half
\begin{frame}[fragile]{}
    \begin{columns}
        \begin{column}{0.5\textwidth}
        \end{column}
        \begin{column}{0.5\textwidth}
        \end{column}
    \end{columns}
\end{frame}

% Slide with bullets
\begin{frame}[fragile]{}
    \begin{itemize}
        \item 
    \end{itemize}
\end{frame}

% Slide with bullets and code
\begin{frame}[fragile]{}
    \begin{itemize}
        \item 
    \end{itemize}
\begin{lstlisting}
Code goes here
\end{lstlisting}
\end{frame}

% Include PDF slides
\includeslides{21-24}{slides-60001/MIT6_0001F16_Lec4_mod.pdf}{Funciones como argumentos}

% Pytutor
\href{URL}{Ver en Python Tutor}


\end{comment}

\end{document}